%!TEX program = pdflatex

\documentclass[12pt,a4paper]{article}

% --- Универсальная настройка движков ---
\usepackage{iftex}

\ifPDFTeX
\usepackage{cmap}
\usepackage[T2A]{fontenc}
\usepackage[utf8]{inputenc}
\usepackage[russian]{babel}
\else
\usepackage{fontspec}
\defaultfontfeatures{Ligatures=TeX}
\usepackage{polyglossia}
\setdefaultlanguage{russian}
\setotherlanguage{english}
\fi

% --- Математические пакеты ---
\usepackage{amsmath,amssymb,amsthm}
\usepackage{icomma}

% --- Верстка, таблицы и графика ---
\usepackage{graphicx}
\usepackage{booktabs}
\usepackage{geometry}
\usepackage{tikz}
\usepackage{xcolor}
\usepackage{microtype}
\usepackage{hyperref}

\geometry{left=2cm,right=2cm,top=2.5cm,bottom=2.5cm}

\hypersetup{
colorlinks=true,
linkcolor=blue!70!black,
citecolor=blue!70!black,
urlcolor=blue!70!black
}

% --- Определение математических операторов ---
\DeclareMathOperator{\Cov}{Cov}

% --- Окружения теорем ---
\newtheorem{theorem}{Теорема}[section]
\newtheorem{definition}{Определение}[section]
\newtheorem{property}{Свойство}[section]
\newtheorem{example}{Пример}[section]
\newtheorem{lemma}{Лемма}[section]

\title{\textbf{Конспект лекций: Теория вероятностей и математическая статистика}\ \large (Оцифрованная и дополненная версия)}
\author{Оцифровано по материалам лекционных тетрадей}
\date{\today}

\begin{document}

\maketitle

\begin{abstract}
Данный документ представляет собой подробную оцифровку лекционных записей по курсу «Теория вероятностей и математическая статистика». В конспект включены ключевые темы, а также подробные алгебраические и комбинаторные доказательства всех основных фактов и теорем. Все доказательства построены на базе дискретных структур и суммирования, что позволяет глубоко понять суть вероятностных процессов без обращения к математическому анализу непрерывных функций.
\end{abstract}

\tableofcontents
\newpage

\section{Элементы теории графов}

\subsection{Основные понятия и определения}
Теория графов является важным инструментом для визуализации и анализа дискретных вероятностных пространств, марковских цепей и комбинаторных задач.

\begin{definition}
\textbf{Граф} $G = (V, E)$ — это математическая модель, состоящая из конечного непустого множества вершин $V$ и множества ребер $E$, представляющего собой связи между парами вершин. Формально, ребро — это неупорядоченная или упорядоченная пара вершин.
\end{definition}

\begin{itemize}
\item \textbf{Изолированная вершина} — вершина со степенью 0 (нет инцидентных ребер).
\item \textbf{Петля} — ребро, связывающее вершину саму с собой.
\item \textbf{Кратные ребра} — несколько ребер, соединяющих одну и ту же пару вершин.
\item \textbf{Простой граф} — граф без петель и кратных ребер. В курсе теории вероятностей чаще всего работают с простыми графами.
\end{itemize}

\subsection{Матричные способы задания графа}
Пусть граф имеет $n = \vert{}V\vert{}$ вершин и $m = \vert{}E\vert{}$ ребер.

\subsubsection{Матрица смежности $B$}
Размер матрицы $B$ составляет $n \times n$. Элемент $B_{ij}$ равен числу ребер, связывающих вершину $v_i$ и вершину $v_j$:
$$B_{ij} = \begin{cases} k, & \text{если вершины } v_i \text{ и } v_j \text{ соединены } k \text{ ребрами}, \\ 0, & \text{в противном случае.} \end{cases}$$
Для неориентированного графа матрица смежности всегда симметрична ($B_{ij} = B_{ji}$).

\subsubsection{Матрица инцидентности $A$}
Размер матрицы $A$ составляет $n \times m$. Для ориентированного графа (орграфа) элементы матрицы определяются следующим образом:
$$A_{ij} = \begin{cases} -1, & \text{если ребро } e_j \text{ выходит из вершины } v_i, \\ 1, & \text{если ребро } e_j \text{ входит в вершину } v_i, \\ 0, & \text{если ребро } e_j \text{ не инцидентно вершине } v_i. \end{cases}$$

\subsection{Связность, мосты и деревья}

\begin{definition}
Ребро $e$ связного графа называется \textbf{мостом}, если при его удалении граф распадается на несколько компонент связности.
\end{definition}

\begin{definition}
\textbf{Деревом} называется связный простой граф, не содержащий циклов. Деревья широко используются при построении деревьев вероятностей.
\end{definition}

\begin{theorem}[О числе ребер в дереве]
Любое дерево с $n \ge 1$ вершинами содержит ровно $m = n - 1$ ребер.
\end{theorem}

\begin{proof}
Доказательство проведем методом математической индукции по числу вершин $n$.
\begin{itemize}
\item \textbf{База индукции:} При $n = 1$ граф состоит из одной вершины. Циклов нет, ребер тоже нет. Формула дает $m = 1 - 1 = 0$. Утверждение верно.
\item \textbf{Индукционный переход:} Предположим, что любое дерево с $k$ вершинами имеет $k - 1$ ребер. Рассмотрим произвольное дерево $T$ с $k + 1$ вершинами. Известно, что в любом дереве с числом вершин больше одной есть хотя бы один «лист» (вершина степени 1).Удалим этот лист и инцидентное ему ребро. Полученный граф $T'$ останется связным (так как удаление листа не разрушает пути между оставшимися вершинами) и в нем не появится циклов (так как мы только удаляли элементы). Следовательно, $T'$ — тоже дерево.

По предположению индукции, так как в $T'$ ровно $k$ вершин, то в нем $k - 1$ ребер. Возвращая удаленную вершину и 1 ребро обратно, получаем, что в исходном дереве $T$ ровно $(k - 1) + 1 = k$ ребер.
\end{itemize}
Утверждение доказано для любого числа вершин $n$.
\end{proof}

\begin{definition}
\textbf{Полный граф} $K_n$ — простой граф с $n$ вершинами, в котором каждые две вершины соединены ребром. Число ребер в $K_n$ равно $m = \frac{n(n-1)}{2}$.
\end{definition}

\subsection{Планарные графы и формула Эйлера}

\begin{definition}
Граф называется \textbf{планарным}, если существует такое его изображение на плоскости, при котором ребра пересекаются только в своих концах (вершинах).
\end{definition}

\begin{theorem}[Формула Эйлера для плоских графов]
Для любого связного планарного графа, уложенного на плоскости, справедливо строгое соотношение:
$$V - E + F = 2,$$
где $V$ — число вершин, $E$ — число ребер, $F$ — число граней (включая одну внешнюю бесконечную грань).
\end{theorem}

\begin{proof}
Доказательство проведем индукцией по числу циклов в графе.
\begin{itemize}
\item \textbf{Случай 1 (без циклов):} Если граф является деревом, то по предыдущей теореме $E = V - 1$. Поскольку циклов нет, граф не ограничивает ни одной внутренней области, следовательно, существует только одна грань — внешняя ($F = 1$).
Проверяем формулу: $V - (V - 1) + 1 = V - V + 1 + 1 = 2$. Верно.
\item \textbf{Случай 2 (с циклами):} Предположим, формула верна для графов с $c$ циклов. Пусть наш граф имеет $c + 1$ циклов. 
Удалим одно любое ребро, принадлежащее какому-либо циклу. 
При этом связность графа не нарушится (так как разорван лишь один обходной путь), но число ребер уменьшится на единицу: $E' = E - 1$. 
Удаление ребра из цикла объединяет две смежные грани (внутреннюю и внешнюю по отношению к удаленному ребру) в одну, следовательно, число граней также уменьшится на 1: $F' = F - 1$. 
Число вершин остается неизменным: $V' = V$.

По индукционному предположению для нового графа выполняется: $V' - E' + F' = 2$.
Подставим исходные переменные: $V - (E - 1) + (F - 1) = V - E + 1 + F - 1 = V - E + F$.
Значение выражения не изменилось и по-прежнему равно 2.
\end{itemize}
Таким образом, формула доказана.
\end{proof}

\begin{theorem}[Теорема Понтрягина — Куратовского]
Граф является планарным тогда и только тогда, когда он не содержит подграфов, гомеоморфных полному графу $K_5$ или полному двудольному графу $K_{3,3}$.
\end{theorem}

\begin{figure}[htbp]
\centering
\begin{tikzpicture}[scale=0.9]
% Полный граф K_5
\begin{scope}[shift={(-4,0)}]
\node at (0,-2.3) {\textbf{а) Полный граф $K_5$}};
\foreach \i in {1,...,5} {
  \node[circle,fill=black,inner sep=2pt] (v\i) at (90-\i*72:1.5) {};
}
\foreach \i in {1,...,5} {
  \foreach \j in {\i,...,5} {
    \draw[thick,blue!70!black] (v\i) -- (v\j);
  }
}
\end{scope}

% Двудольный граф K_{3,3}
\begin{scope}[shift={(4,0)}]
\node at (0,-2.3) {\textbf{б) Двудольный граф $K_{3,3}$}};
\foreach \i in {1,2,3} {
  \node[circle,fill=black,inner sep=2pt] (a\i) at (-1.3, 2-\i*1) {};
  \node[circle,fill=black,inner sep=2pt] (b\i) at (1.3, 2-\i*1) {};
}
\foreach \i in {1,2,3} {
  \foreach \j in {1,2,3} {
    \draw[thick,red!70!black] (a\i) -- (b\j);
  }
}
\end{scope}
\end{tikzpicture}
\caption{Классические непланарные графы (критерий Куратовского)}
\end{figure}
\section{Случайные события и вероятностное пространство}\subsection{Пространство элементарных исходов}Основной концепцией теории вероятностей является формализация случайного эксперимента.
\textbf{Пространство элементарных исходов} $\Omega$ — это множество всех возможных взаимно исключающих исходов $\omega$, которыми может завершиться эксперимент.Событие $A$ — это любое подмножество исходов $A \subseteq \Omega$.\subsection{Классическое определение вероятности}Если эксперимент удовлетворяет двум условиям:Число всех исходов $\vert{}\Omega\vert{} = n$ конечно.Все исходы равновозможны (симметрия эксперимента).То вероятность наступления события $A$, состоящего из $m$ исходов, определяется как отношение:
$$P(A) = \frac{\vert{}A\vert{}}{\vert{}\Omega\vert{}} = \frac{m}{n}.$$\subsection{Теоремы сложения вероятностей}\begin{theorem}[Теорема сложения для несовместных событий]
Если события $A$ и $B$ не могут наступить одновременно (т.е. их пересечение пусто: $A \cap B = \varnothing$), то вероятность того, что наступит хотя бы одно из них, равна сумме их вероятностей:
$$P(A \cup B) = P(A) + P(B).$$
\end{theorem}\begin{proof}
Пусть общее число исходов равно $n$. Событию $A$ благоприятствует $m_A$ исходов, событию $B$ — $m_B$ исходов. Поскольку $A$ и $B$ не пересекаются, ни один исход не благоприятствует обоим событиям одновременно. Значит, событию $A \cup B$ благоприятствует ровно $m_A + m_B$ исходов.
$$P(A \cup B) = \frac{m_A + m_B}{n} = \frac{m_A}{n} + \frac{m_B}{n} = P(A) + P(B).$$
\end{proof}\begin{theorem}[Общая теорема сложения вероятностей]
Для любых двух событий $A$ и $B$ (совместных или несовместных):
$$P(A \cup B) = P(A) + P(B) - P(A \cap B).$$
\end{theorem}\begin{proof}
Представим событие $A \cup B$ как объединение двух несовместных событий: самого события $A$ и той части события $B$, которая не входит в $A$. Формально: $A \cup B = A \cup (B \setminus A)$.
Так как эти события несовместны, применим предыдущую теорему:
$$P(A \cup B) = P(A) + P(B \setminus A). \quad (1)$$
В свою очередь, само событие $B$ можно разбить на две несовместные части: те исходы, которые входят только в $B$, и те, которые входят и в $B$, и в $A$.
$$B = (B \setminus A) \cup (A \cap B).$$
Вероятность $B$ равна сумме их вероятностей:
$$P(B) = P(B \setminus A) + P(A \cap B) \implies P(B \setminus A) = P(B) - P(A \cap B).$$
Подставим полученное выражение для $P(B \setminus A)$ в уравнение (1):
$$P(A \cup B) = P(A) + P(B) - P(A \cap B).$$
Теорема доказана.
\end{proof}\begin{theorem}[Формула включений-исключений для 3 событий]
$$P(A \cup B \cup C) = P(A) + P(B) + P(C) - P(A \cap B) - P(A \cap C) - P(B \cap C) + P(A \cap B \cap C).$$
\end{theorem}\begin{proof}
Обозначим $D = A \cup B$. Тогда $A \cup B \cup C = D \cup C$. Применим общую теорему сложения для двух событий $D$ и $C$:
$$P(D \cup C) = P(D) + P(C) - P(D \cap C).$$
Раскроем $P(D)$: $P(D) = P(A \cup B) = P(A) + P(B) - P(A \cap B)$. \
Теперь раскроем $P(D \cap C)$:
$$D \cap C = (A \cup B) \cap C = (A \cap C) \cup (B \cap C).$$
Вероятность этого объединения (по теореме сложения для двух событий):
$$P((A \cap C) \cup (B \cap C)) = P(A \cap C) + P(B \cap C) - P((A \cap C) \cap (B \cap C)).$$
Заметим, что $(A \cap C) \cap (B \cap C) = A \cap B \cap C$. \
Собираем все воедино:
\begin{align*}
P(A \cup B \cup C) &= [P(A) + P(B) - P(A \cap B)] + P(C) - [P(A \cap C) + P(B \cap C) - P(A \cap B \cap C)] \
&= P(A) + P(B) + P(C) - P(A \cap B) - P(A \cap C) - P(B \cap C) + P(A \cap B \cap C).
\end{align*}
\end{proof}
\section{Геометрическая и условная вероятность}

\subsection{Геометрическая вероятность}
Классическое определение вероятности неприменимо, если число возможных исходов бесконечно. Для таких случаев вводится геометрическое определение.

Пусть пространство исходов~$\Omega$ представляет собой некоторую область (отрезок, плоскую фигуру, пространственное тело), имеющую конечную меру $\mu(\Omega)$ (длину, площадь, объем соответственно). Если вероятность попадания случайной точки в любую часть области пропорциональна мере этой части и не зависит от ее расположения и формы, то вероятность события $A \subseteq \Omega$ равна:
\[
\Prob(A) = \frac{\mu(A)}{\mu(\Omega)}.
\]

\begin{example}[Задача о спагетти]
Спагетти длиной~1 разрывается в двух случайных точках $x$ и $y$ ($0 \le x \le y \le 1$). Какова вероятность того, что из трех образовавшихся отрезков длинами $x$, $y-x$ и $1-y$ можно сложить треугольник?

\textbf{Решение:} По неравенству треугольника, длина ни одного из отрезков не должна превышать $0{,}5$:
\[
\begin{cases} x \le 0{,}5, \\ y - x \le 0{,}5 \implies y \le x + 0{,}5, \\ 1 - y \le 0{,}5 \implies y \ge 0{,}5. \end{cases}
\]
Область всех возможных исходов $\Omega = \{(x,y): 0 \le x \le y \le 1\}$ имеет площадь $S_\Omega = 0{,}5$. Область благоприятных исходов~$A$ образует треугольник с площадью $S_A = 0{,}125$. Искомая вероятность:
\[
\Prob(A) = \frac{S_A}{S_\Omega} = \frac{0{,}125}{0{,}5} = \frac{1}{4} = 0{,}25.
\]
\end{example}

\begin{example}[Задача о корнях квадратного уравнения]
Коэффициенты $p$ и $q$ выбираются наудачу из отрезка $[-1, 1]$. Какова вероятность того, что квадратное уравнение $x^2 + px + q = 0$ имеет действительные корни?

\textbf{Решение:} Уравнение имеет действительные корни при $D = p^2 - 4q \ge 0 \implies q \le \frac{p^2}{4}$. 
Область $\Omega = [-1,1] \times [-1,1]$, $S_\Omega = 4$. 
Площадь области $A = \{(p,q) \in [-1,1]^2 : q \le \frac{p^2}{4}\}$ вычисляется через интеграл:
\[
S_A = 2 + \int_{-1}^1 \frac{p^2}{4} \, dp = 2 + \frac{1}{6} = \frac{13}{6} \approx 2{,}167.
\]
Вероятность $\Prob(A) = \frac{13/6}{4} = \frac{13}{24} \approx 0{,}5417$.
\end{example}

\subsection{Условная вероятность и теорема умножения}

\begin{definition}
\textbf{Условной вероятностью} $\Prob(A \mid B)$ события~$A$ при условии, что произошло событие~$B$ (где $\Prob(B) > 0$), называется вероятность наступления~$A$ в суженном пространстве исходов, где достоверным стало событие~$B$.
Определяется она формулой:
\[
\Prob(A \mid B) = \frac{\Prob(A \cap B)}{\Prob(B)}.
\]
\end{definition}

Из этого определения напрямую вытекает важная теорема.

\begin{theorem}[Теорема умножения вероятностей]
Вероятность совместного наступления двух событий~$A$ и~$B$ равна произведению вероятности одного из них на условную вероятность другого:
\[
\Prob(A \cap B) = \Prob(B) \cdot \Prob(A \mid B) = \Prob(A) \cdot \Prob(B \mid A).
\]
\end{theorem}

\begin{proof}
В классической вероятностной модели пусть $n$~--- общее число исходов. $m_B$~--- число исходов, благоприятствующих~$B$. $m_{AB}$~--- число исходов, благоприятствующих совместному наступлению $A$ и $B$ ($A \cap B$).
Тогда, если известно, что событие~$B$ произошло, пространство исходов сужается до $m_B$ вариантов. Из них событию~$A$ благоприятствует $m_{AB}$ исходов.
По классическому определению для суженного пространства:
\[
\Prob(A \mid B) = \frac{m_{AB}}{m_B}.
\]
Перепишем дробь, разделив числитель и знаменатель на общее число исходов~$n$:
\[
\Prob(A \mid B) = \frac{m_{AB} / n}{m_B / n} = \frac{\Prob(A \cap B)}{\Prob(B)}.
\]
Умножив обе части на $\Prob(B)$, получаем $\Prob(A \cap B) = \Prob(B) \cdot \Prob(A \mid B)$. Меняя роли $A$ и $B$, аналогично выводится $\Prob(A \cap B) = \Prob(A) \cdot \Prob(B \mid A)$.
\end{proof}

\subsection{Независимость событий}

\begin{definition}
События~$A$ и~$B$ называются \textbf{независимыми}, если наступление одного не меняет вероятности наступления другого, то есть:
\[
\Prob(A \mid B) = \Prob(A).
\]
\end{definition}

Подставляя это в теорему умножения, получаем строгий алгебраический критерий независимости двух событий:
\[
\Prob(A \cap B) = \Prob(A) \cdot \Prob(B).
\]

\subsection{Парадокс Монти Холла}

Парадокс Монти Холла~--- одна из самых известных задач теории вероятностей, наглядно иллюстрирующая контринтуитивную природу условной вероятности.

\begin{example}[Парадокс Монти Холла]
В телеигре участнику предлагают выбрать одну из трех закрытых дверей ($1, 2, 3$). За одной из дверей находится автомобиль, за двумя другими~--- козлы. 
\begin{enumerate}
\item Игрок выбирает одну из дверей (пусть для определенности дверь~1).
\item Ведущий (Монти Холл), который знает, где именно спрятан автомобиль, открывает одну из двух оставшихся дверей (например, дверь~3), за которой гарантированно оказывается козел.
\item Ведущий предлагает игроку выбор: остаться при своем первоначальном решении (дверь~1) или сменить дверь на оставшуюся закрытую (дверь~2).
\end{enumerate}
Выгодно ли сменить свой выбор?

\textbf{Строгое решение через условную вероятность:} \\
Обозначим через $C_i$ гипотезу о том, что автомобиль находится за дверью~$i$ ($i \in \{1, 2, 3\}$). До начала игры все положения равновероятны:
\[
\Prob(C_1) = \Prob(C_2) = \Prob(C_3) = \frac{1}{3}.
\]
Пусть игрок выбрал дверь~1. Обозначим через $M_3$ событие «ведущий открыл дверь~3». 
Определим условные вероятности действий ведущего при различных положениях автомобиля:
\begin{itemize}
\item Если приз за дверью~1 ($C_1$), у ведущего есть выбор между дверями~2 и~3 (за обеими козлы). Он выбирает наудачу: $\Prob(M_3 \mid C_1) = \frac{1}{2}$.
\item Если приз за дверью~2 ($C_2$), ведущий не может открыть выбранную игроком дверь~1 и не может открыть выигрышную дверь~2. У него нет выбора: $\Prob(M_3 \mid C_2) = 1$.
\item Если приз за дверью~3 ($C_3$), ведущий никогда не откроет дверь с автомобилем: $\Prob(M_3 \mid C_3) = 0$.
\end{itemize}

По формуле полной вероятности найдем безусловную вероятность того, что ведущий откроет именно дверь~3:
\[
\Prob(M_3) = \sum_{i=1}^3 \Prob(C_i)\Prob(M_3 \mid C_i) = \frac{1}{3}\cdot\frac{1}{2} + \frac{1}{3}\cdot 1 + \frac{1}{3}\cdot 0 = \frac{1}{6} + \frac{1}{3} = \frac{1}{2}.
\]

Теперь вычислим апостериорные (условные) вероятности выигрыша:
\begin{align*}
\Prob(C_1 \mid M_3) &= \frac{\Prob(C_1 \cap M_3)}{\Prob(M_3)} = \frac{\frac{1}{3} \cdot \frac{1}{2}}{\frac{1}{2}} = \frac{1}{3} \quad \text{(шанс победы без смены двери)}, \\
\Prob(C_2 \mid M_3) &= \frac{\Prob(C_2 \cap M_3)}{\Prob(M_3)} = \frac{\frac{1}{3} \cdot 1}{\frac{1}{2}} = \frac{2}{3} \quad \text{(шанс победы при смене двери)}.
\end{align*}

Следовательно, смена двери удваивает вероятность выигрыша автомобиля с $\frac{1}{3}$ до $\frac{2}{3}$.
\end{example}

\begin{figure}[htbp]
\centering
\begin{tikzpicture}[
  node distance=1.5cm and 2.5cm,
  box/.style={rectangle, draw=blue!70!black, thick, fill=blue!5, rounded corners=3pt, align=center, inner sep=6pt, font=\small},
  action/.style={rectangle, draw=gray!70!black, fill=gray!10, rounded corners=2pt, align=center, inner sep=4pt, font=\footnotesize},
  win/.style={rectangle, draw=teal!70!black, thick, fill=teal!15, rounded corners=2pt, font=\small\bfseries},
  lose/.style={rectangle, draw=red!70!black, thick, fill=red!15, rounded corners=2pt, font=\small\bfseries},
  arrow/.style={-{Stealth[length=2.5mm]}, thick, gray!60!black}
]
  \node[box] (start) {Игрок выбирает дверь~1};

  \node[action, below left=1.2cm and 2.2cm of start] (c1) {Авто за дверью 1\\($p = 1/3$)};
  \node[action, below=1.2cm of start] (c2) {Авто за дверью 2\\($p = 1/3$)};
  \node[action, below right=1.2cm and 2.2cm of start] (c3) {Авто за дверью 3\\($p = 1/3$)};

  \draw[arrow] (start) -- (c1);
  \draw[arrow] (start) -- (c2);
  \draw[arrow] (start) -- (c3);

  \node[action, below=0.8cm of c1] (host1) {Ведущий открывает\\дверь 2 или 3};
  \node[action, below=0.8cm of c2] (host2) {Ведущий обязан\\открыть дверь 3};
  \node[action, below=0.8cm of c3] (host3) {Ведущий обязан\\открыть дверь 2};

  \draw[arrow] (c1) -- (host1);
  \draw[arrow] (c2) -- (host2);
  \draw[arrow] (c3) -- (host3);

  \node[lose, below=0.8cm of host1] (res1) {Смена $\to$ Проигрыш};
  \node[win, below=0.8cm of host2] (res2) {Смена $\to$ Автомобиль!};
  \node[win, below=0.8cm of host3] (res3) {Смена $\to$ Автомобиль!};

  \draw[arrow] (host1) -- (res1);
  \draw[arrow] (host2) -- (res2);
  \draw[arrow] (host3) -- (res3);
\end{tikzpicture}
\caption{Дерево исходов при стратегии обязательной смены двери в парадоксе Монти Холла}
\label{fig:monty_hall}
\end{figure}

\textbf{Интуитивное объяснение:} В момент первоначального выбора игрок ошибается с вероятностью $\frac{2}{3}$. Если игрок выбрал козла (что происходит в $2$ случаях из $3$), у ведущего нет свободы маневра: он обязан открыть единственную оставшуюся дверь с другим козлом. Тем самым ведущий непреднамеренно указывает игроку пальцем на дверь с автомобилем. Поэтому смена выбора выигрывает ровно тогда, когда исходный выбор был ошибочным, то есть с вероятностью $\frac{2}{3}$.
\section{Формула полной вероятности и формула Байеса}

Во многих задачах случайный эксперимент протекает в два этапа. На первом этапе реализуется одно из нескольких возможных условий (гипотез), а на втором — происходит интересующее нас событие.

\subsection{Полная группа гипотез}

\begin{definition}
События $H_1, H_2, \dots, H_n$ образуют \textbf{полную группу гипотез}, если выполняются три условия:
\begin{enumerate}
\item Они попарно несовместны: $H_i \cap H_j = \varnothing$ для любых $i \neq j$.
\item В результате эксперимента обязательно произойдет хотя бы одно из них: $\bigcup_{i=1}^n H_i = \Omega$.
\item Сумма их вероятностей равна единице: $\sum_{i=1}^n P(H_i) = 1$.
\end{enumerate}
\end{definition}

\subsection{Формула полной вероятности}

\begin{theorem}[Формула полной вероятности]
Вероятность события $A$, которое может наступить лишь при условии появления одного из несовместных событий $H_1, \dots, H_n$, образующих полную группу, равна сумме произведений вероятностей каждой из гипотез на соответствующую условную вероятность события $A$:
$$P(A) = \sum_{i=1}^n P(H_i) \cdot P(A\vert{}H_i).$$
\end{theorem}

\begin{proof}
Так как гипотезы $H_1, \dots, H_n$ образуют полное пространство $\Omega$, событие $A$ можно представить как его пересечение со всем пространством:
$$A = A \cap \Omega = (A \cap H_1) \cup (A \cap H_2) \cup \dots \cup (A \cap H_n).$$
События $(A \cap H_i)$ попарно несовместны. Применим аксиому сложения вероятностей:
$$P(A) = \sum_{i=1}^n P(A \cap H_i).$$
Согласно теореме умножения вероятностей $P(A \cap H_i) = P(H_i) \cdot P(A\vert{}H_i)$. Подставляем в сумму и получаем искомую формулу.
\end{proof}

\subsection{Формула Байеса (Теорема гипотез)}

\begin{theorem}[Формула Байеса]
Если событие $A$ уже произошло ($P(A) > 0$), то вероятность того, что оно было вызвано именно гипотезой $H_k$, пересчитывается (уточняется) по формуле:
$$P(H_k\vert{}A) = \frac{P(H_k) \cdot P(A\vert{}H_k)}{P(A)} = \frac{P(H_k) \cdot P(A\vert{}H_k)}{\sum_{i=1}^n P(H_i) \cdot P(A\vert{}H_i)}.$$
\end{theorem}

\begin{proof}
Запишем вероятность совместного наступления гипотезы $H_k$ и события $A$ двумя способами:
$$P(H_k \cap A) = P(A) \cdot P(H_k\vert{}A) \quad \text{и} \quad P(H_k \cap A) = P(H_k) \cdot P(A\vert{}H_k).$$
Приравняем правые части и разделим на $P(A)$:
$$P(H_k\vert{}A) = \frac{P(H_k) \cdot P(A\vert{}H_k)}{P(A)}.$$
\end{proof}

\subsection{Визуализация: Деревья вероятностей}

Для наглядного решения задач на формулу полной вероятности и формулу Байеса удобно использовать \textbf{деревья вероятностей}.

\begin{example}[Диагностика спама]
Пусть $8\%$ сообщений являются спамом (гипотеза $H_1$), а $92\%$ — нормальными письмами (гипотеза $H_2$). Вероятность того, что фильтр пометит спам как спам (событие $A$), равна $0{,}95$. Вероятность ложного срабатывания на нормальном письме равна $0{,}01$. 

Найдем общую вероятность срабатывания фильтра $P(A)$ и вероятность того, что помеченное письмо — действительно спам $P(H_1|A)$.

\begin{center}
\begin{tikzpicture}[
  level 1/.style={sibling distance=6cm, level distance=2.5cm},
  level 2/.style={sibling distance=3cm, level distance=2.5cm},
  every node/.style={align=center, fill=white, inner sep=2pt}
]
\node {Поступающее\\письмо}
  child { node {$H_1$ (Спам)}
    child { node {$A$ (Метка)} edge from parent node[left=2pt] {$P(A|H_1) = 0{,}95$} }
    child { node {$\overline{A}$ (Нет)} edge from parent node[right=2pt] {$0{,}05$} }
    edge from parent node[left=5pt, fill=none] {$P(H_1) = 0{,}08$}
  }
  child { node {$H_2$ (Норма)}
    child { node {$A$ (Метка)} edge from parent node[left=2pt] {$P(A|H_2) = 0{,}01$} }
    child { node {$\overline{A}$ (Нет)} edge from parent node[right=2pt] {$0{,}99$} }
    edge from parent node[right=5pt, fill=none] {$P(H_2) = 0{,}92$}
  };
\end{tikzpicture}
\end{center}

По формуле полной вероятности суммируем все пути, ведущие к событию $A$:
$$P(A) = 0{,}08 \cdot 0{,}95 + 0{,}92 \cdot 0{,}01 = 0{,}076 + 0{,}0092 = 0{,}0852.$$

По формуле Байеса вычисляем долю «истинного спама» среди всех сработавших меток:
$$P(H_1|A) = \frac{P(H_1)P(A|H_1)}{P(A)} = \frac{0{,}076}{0{,}0852} \approx 0{,}892 \quad (89{,}2\%).$$
\end{example}
\section{Дискретные случайные величины}

\subsection{Определение и закон распределения}

\begin{definition}
\textbf{Случайной величиной} $X$ называется числовая функция, ставящая в соответствие каждому элементарному исходу $\omega \in \Omega$ некоторое действительное число $X(\omega)$.
\end{definition}

Случайная величина называется \textbf{дискретной}, если множество ее возможных значений конечно или счетно. Закон ее распределения задается таблицей (рядом распределения), где перечислены значения $x_i$ и соответствующие им вероятности $p_i = P(X = x_i)$. Сумма всех вероятностей обязана равняться единице: $\sum p_i = 1$.

\subsection{Математическое ожидание и его свойства}

\begin{definition}
\textbf{Математическое ожидание} $E(X)$ дискретной случайной величины — это сумма произведений всех ее возможных значений на их вероятности:
$$E(X) = \sum_{i} x_i p_i.$$
Оно характеризует «центр тяжести» распределения.
\end{definition}

\begin{theorem}[Свойства математического ожидания]
Пусть $X$ и $Y$ — случайные величины, а $c, a, b$ — константы.
\begin{enumerate}
\item $E(c) = c$.
\item $E(aX + b) = a E(X) + b$.
\item $E(XY) = E(X) E(Y)$ (только если $X$ и $Y$ независимы).
\end{enumerate}
\end{theorem}

\begin{proof}
Докажем свойства, опираясь исключительно на определение суммы:

1. Константу $c$ можно рассматривать как дискретную случайную величину, которая принимает единственное значение $c$ с вероятностью $p_1 = 1$. Тогда по определению:
$$E(c) = c \cdot 1 = c.$$

2. Рассмотрим новую случайную величину $Z = aX + b$. Ее значения равны $(a x_i + b)$, и они принимаются с теми же вероятностями $p_i$.
$$E(aX + b) = \sum_{i} (a x_i + b) p_i = \sum_{i} (a x_i p_i + b p_i).$$
Разбиваем сумму на две:
$$\sum_{i} a x_i p_i + \sum_{i} b p_i = a \left( \sum_{i} x_i p_i \right) + b \left( \sum_{i} p_i \right).$$
Так как $\sum x_i p_i = E(X)$, а $\sum p_i = 1$ (условие нормировки), получаем $a E(X) + b \cdot 1 = a E(X) + b$.

3. Доказательство для произведения независимых величин приведено в Главе 6 (Двумерные случайные величины), где формализуется совместное распределение.
\end{proof}

\subsection{Дисперсия и ее свойства}

\begin{definition}
\textbf{Дисперсия} $D(X)$ — это математическое ожидание квадрата отклонения случайной величины от ее математического ожидания:
$$D(X) = E\left((X - E(X))^2\right).$$
\end{definition}

\begin{theorem}[Вычислительная формула дисперсии]
Дисперсия равна разности математического ожидания квадрата величины и квадрата ее математического ожидания: $D(X) = E(X^2) - (E(X))^2$.
\end{theorem}

\begin{proof}
Обозначим для краткости константу $\mu = E(X)$. Тогда:
$$D(X) = E\left((X - \mu)^2\right).$$
Раскроем квадрат под знаком математического ожидания алгебраически: $(X - \mu)^2 = X^2 - 2\mu X + \mu^2$.
Используя свойство линейности математического ожидания, разобьем выражение:
$$D(X) = E(X^2 - 2\mu X + \mu^2) = E(X^2) - E(2\mu X) + E(\mu^2).$$
Поскольку $\mu$ — это число (константа), то $E(2\mu X) = 2\mu E(X)$ и $E(\mu^2) = \mu^2$.
Подставляем обратно $E(X) = \mu$:
$$D(X) = E(X^2) - 2\mu \cdot \mu + \mu^2 = E(X^2) - 2\mu^2 + \mu^2 = E(X^2) - \mu^2 = E(X^2) - (E(X))^2.$$
Теорема доказана.
\end{proof}

\begin{theorem}[Свойства дисперсии]
\begin{enumerate}
\item $D(c) = 0$.
\item $D(X + c) = D(X)$.
\item $D(aX) = a^2 D(X)$.
\end{enumerate}
\end{theorem}

\begin{proof}
\begin{enumerate}
\item $D(c) = E(c^2) - (E(c))^2 = c^2 - c^2 = 0$. 
\item $D(X + c) = E(((X+c) - E(X+c))^2) = E((X + c - E(X) - c)^2) = E((X - E(X))^2) = D(X)$. Сдвиг распределения не меняет его разброс. 
\item $D(aX) = E((aX - E(aX))^2) = E((aX - aE(X))^2) = E(a^2 (X - E(X))^2)$. Выносим константу $a^2$ за знак $E$: получаем $a^2 E((X - E(X))^2) = a^2 D(X)$.
\end{enumerate}
\end{proof}
\section{Двумерные случайные величины и корреляция}

\subsection{Совместное распределение и математическое ожидание суммы}

Многие эксперименты требуют изучения сразу двух параметров $X$ и $Y$. Совместное распределение пары дискретных случайных величин задается матрицей вероятностей $p_{ij} = P(X = x_i, Y = y_j)$.

\begin{theorem}[Линейность математического ожидания для суммы]
Для любых двух случайных величин $X$ и $Y$ (независимо от того, зависимы они или нет) выполняется:
$$E(X + Y) = E(X) + E(Y).$$
\end{theorem}

\begin{proof}
По определению математического ожидания функции от двумерной случайной величины, мы должны просуммировать значения суммы $(x_i + y_j)$ по всем возможным исходам с их вероятностями $p_{ij}$:
\begin{align*}
E(X + Y) &= \sum_{i} \sum_{j} (x_i + y_j) p_{ij} = \sum_{i} \sum_{j} (x_i p_{ij} + y_j p_{ij}) \\
&= \sum_{i} \sum_{j} x_i p_{ij} + \sum_{i} \sum_{j} y_j p_{ij}.
\end{align*}
В первой двойной сумме вынесем $x_i$ (оно не зависит от индекса $j$), а во второй вынесем $y_j$ (не зависит от $i$):
$$E(X + Y) = \sum_{i} x_i \left( \sum_{j} p_{ij} \right) + \sum_{j} y_j \left( \sum_{i} p_{ij} \right).$$
Заметим, что суммирование совместной вероятности $p_{ij}$ по всем $j$ дает маргинальную (одномерную) вероятность $P(X = x_i)$. Аналогично, сумма по $i$ дает $P(Y = y_j)$:
$$E(X + Y) = \sum_{i} x_i P(X = x_i) + \sum_{j} y_j P(Y = y_j) = E(X) + E(Y).$$
Теорема доказана.
\end{proof}

\begin{theorem}[Математическое ожидание произведения независимых СВ]
Если $X$ и $Y$ независимы, то $E(XY) = E(X)E(Y)$.
\end{theorem}

\begin{proof}
Если $X$ и $Y$ независимы, то вероятность совместного наступления равна произведению вероятностей: $p_{ij} = P(X = x_i) \cdot P(Y = y_j)$.
$$E(XY) = \sum_i \sum_j (x_i y_j) p_{ij} = \sum_i \sum_j (x_i y_j) P(X = x_i) P(Y = y_j).$$
Сгруппируем члены, зависящие от $i$ и от $j$:
$$E(XY) = \left( \sum_i x_i P(X = x_i) \right) \cdot \left( \sum_j y_j P(Y = y_j) \right) = E(X) \cdot E(Y).$$
\end{proof}

\subsection{Ковариация и корреляция}

\begin{definition}
\textbf{Ковариацией} случайных величин $X$ и $Y$ называется математическое ожидание произведения их отклонений от своих математических ожиданий:
$$\Cov(X, Y) = E\left((X - E(X))(Y - E(Y))\right).$$
\end{definition}

\begin{theorem}[Вычислительная формула и независимость]
\begin{enumerate}
\item $\Cov(X, Y) = E(XY) - E(X)E(Y)$.
\item Если $X$ и $Y$ независимы, то $\Cov(X, Y) = 0$.
\end{enumerate}
\end{theorem}

\begin{proof}
\begin{enumerate}
\item Пусть $\mu_X = E(X)$ и $\mu_Y = E(Y)$. Раскроем скобки:
\begin{align*}
\Cov(X,Y) &= E((X - \mu_X)(Y - \mu_Y)) = E(XY - X\mu_Y - Y\mu_X + \mu_X\mu_Y) \\
&= E(XY) - \mu_Y E(X) - \mu_X E(Y) + \mu_X\mu_Y \\
&= E(XY) - \mu_Y\mu_X - \mu_X\mu_Y + \mu_X\mu_Y = E(XY) - \mu_X\mu_Y = E(XY) - E(X)E(Y).
\end{align*}

\item Если $X$ и $Y$ независимы, то по доказанной выше теореме $E(XY) = E(X)E(Y)$.
Подставляем это в формулу: $\Cov(X,Y) = E(X)E(Y) - E(X)E(Y) = 0$.
\end{enumerate}
\end{proof}

\begin{theorem}[Дисперсия суммы]
$D(X + Y) = D(X) + D(Y) + 2\,\Cov(X, Y)$.
\end{theorem}

\begin{proof}
По определению дисперсии: $D(X+Y) = E\left( ((X+Y) - E(X+Y))^2 \right)$.
Перегруппируем слагаемые:
$$D(X+Y) = E\left( ((X - E(X)) + (Y - E(Y)))^2 \right).$$
Возведем в квадрат по формуле $(a+b)^2 = a^2 + 2ab + b^2$:
\begin{align*}
D(X+Y) &= E\left( (X - E(X))^2 + 2(X - E(X))(Y - E(Y)) + (Y - E(Y))^2 \right) \\
&= E\left((X - E(X))^2\right) + 2E\left((X - E(X))(Y - E(Y))\right) + E\left((Y - E(Y))^2\right) \\
&= D(X) + 2\Cov(X, Y) + D(Y).
\end{align*}
\end{proof}
\section{Предельные теоремы теории вероятностей}

Закон больших чисел (ЗБЧ) математически обосновывает тот факт, что при большом числе случайных экспериментов их средний результат перестает быть случайным и стремится к определенному постоянному значению. Для доказательства ЗБЧ используются неравенства Маркова и Чебышева.

\subsection{Неравенство Маркова}
\begin{theorem}[Неравенство Маркова]
Пусть $X \ge 0$ — неотрицательная дискретная случайная величина с конечным математическим ожиданием $E(X)$. Тогда для любого положительного числа $\varepsilon > 0$:
$$P(X \ge \varepsilon) \le \frac{E(X)}{\varepsilon}.$$
\end{theorem}

\begin{proof}
Запишем определение математического ожидания для дискретной случайной величины. Так как $X \ge 0$, все ее значения $x_i \ge 0$.
Разделим сумму на две части: значения, которые меньше $\varepsilon$, и значения, которые больше или равны $\varepsilon$.
$$E(X) = \sum_{x_i < \varepsilon} x_i P(X = x_i) + \sum_{x_i \ge \varepsilon} x_i P(X = x_i).$$
Первая сумма неотрицательна. Если мы ее отбросим, мы только уменьшим общее выражение:
$$E(X) \ge \sum_{x_i \ge \varepsilon} x_i P(X = x_i).$$
В оставшейся сумме каждый элемент $x_i \ge \varepsilon$. Если мы заменим каждый $x_i$ на меньшее число $\varepsilon$, сумма снова уменьшится:
$$E(X) \ge \varepsilon \sum_{x_i \ge \varepsilon} P(X = x_i) = \varepsilon \cdot P(X \ge \varepsilon).$$
Разделив обе части на $\varepsilon > 0$, получаем искомое неравенство: $P(X \ge \varepsilon) \le \frac{E(X)}{\varepsilon}$.
\end{proof}

\subsection{Неравенство Чебышева}
\begin{theorem}[Неравенство Чебышева]
Для любой случайной величины $X$ с конечными математическим ожиданием $\mu = E(X)$ и дисперсией $D(X)$, для любого $\varepsilon > 0$ справедливо:
$$P(\vert{}X - \mu\vert{} \ge \varepsilon) \le \frac{D(X)}{\varepsilon^2}.$$
\end{theorem}

\begin{proof}
Рассмотрим новую случайную величину $Y = (X - \mu)^2$. Эта величина неотрицательна, а ее математическое ожидание по определению равно дисперсии: $E(Y) = D(X)$.
Событие $\vert{}X - \mu\vert{} \ge \varepsilon$ эквивалентно событию $(X - \mu)^2 \ge \varepsilon^2$, то есть $Y \ge \varepsilon^2$.
Применим доказанное ранее неравенство Маркова для величины $Y$ и константы $\varepsilon^2$:
$$P(Y \ge \varepsilon^2) \le \frac{E(Y)}{\varepsilon^2} \implies P(\vert{}X - \mu\vert{} \ge \varepsilon) \le \frac{D(X)}{\varepsilon^2}.$$
\end{proof}

\subsection{Схема Бернулли и Биномиальное распределение}
Число успехов $S_n$ в $n$ независимых испытаниях Бернулли подчиняется биномиальному закону с вероятностью:
$$P(S_n = k) = C_n^k p^k q^{n-k}.$$
Ее основные числовые характеристики (математическое ожидание и дисперсия) выражаются формулами:
$$E(S_n) = n p, \quad D(S_n) = n p q.$$

\subsection{Законы больших чисел}

\begin{theorem}[Закон больших чисел в форме Бернулли]
Относительная частота $\frac{S_n}{n}$ сходится по вероятности к вероятности успеха $p$:
$$P\left(\left\vert\frac{S_n}{n} - p\right\vert \ge \varepsilon\right) \le \frac{p q}{n \varepsilon^2} \xrightarrow[n \to \infty]{} 0.$$
\end{theorem}

\begin{theorem}[Теорема Чебышева]
Пусть $X_1, X_2, \dots, X_n$ — последовательность попарно независимых случайных величин с математическими ожиданиями $E(X_i) = \mu_i$ и ограниченными дисперсиями $D(X_i) \le C$. Тогда среднее арифметическое этих величин сходится по вероятности к среднему их математических ожиданий:
$$\lim_{n \to \infty} P\left(\left\vert{}\frac{1}{n}\sum_{i=1}^n X_i - \frac{1}{n}\sum_{i=1}^n \mu_i\right\vert{} \ge \varepsilon\right) = 0.$$
\end{theorem}

\begin{proof}
Рассмотрим случайную величину $Z_n = \frac{1}{n}\sum_{i=1}^n X_i$.
По свойствам линейности (см. Теорему \ref{thm:ev_properties} на стр. \pageref{thm:ev_properties}), ее математическое ожидание равно $E(Z_n) = \frac{1}{n}\sum_{i=1}^n \mu_i$.
Поскольку $X_i$ независимы, дисперсия суммы равна сумме дисперсий (согласно Теореме \ref{thm:var_properties}):
$$D(Z_n) = D\left(\frac{1}{n}\sum_{i=1}^n X_i\right) = \frac{1}{n^2} \sum_{i=1}^n D(X_i).$$
Так как каждая дисперсия ограничена константой $C$, то:
$$D(Z_n) \le \frac{1}{n^2} \sum_{i=1}^n C = \frac{1}{n^2} (n \cdot C) = \frac{C}{n}.$$
Применим неравенство Чебышева к величине $Z_n$:
$$P(\vert{}Z_n - E(Z_n)\vert{} \ge \varepsilon) \le \frac{D(Z_n)}{\varepsilon^2} \le \frac{C}{n \varepsilon^2}.$$
При стремлении $n \to \infty$ правая часть стремится к нулю. Следовательно, вероятность отклонения от математического ожидания тоже стремится к нулю.
\end{proof}

\end{document}